\documentclass[12pt]{article}

\begin{document}

PART I
BASIC
CONCEPTS

\hfill

In the first three chapters we present an overview
and the basic concepts, beginning with critical
points in the one-dimensional case and proceeding
to the corresponding idea of critical curves in two
dimensions.

\hfill

\textbf{CHAPTER I
INTRODUCTION}

\hfill

\textbf{1.1 BACKGROUND}

\hfill

Our goal is to present the fundamentals of two-dimensional
(2D) iteration theory through examples, with extensive graphics
(for which the 2D context is ideal) and few mathematical symbols.1

\hfill

We illustrate all the basic ideas with hand drawings and monochrome
computer graphics in the book, and again with movies (fullmotion
video animations in color) on the companion CD-ROM.
We do not assume a knowledge of higher mathematics. But we
do acknowledge that our subject is a branch of pure mathematics,
and a deeper understanding requires some topology and geometry.
A hint of this is presented in the appendices, where a more rigorous
approach is introduced.

\hfill

\textbf{1.2 HISTORY}

\hfill

The study of chaos in ID iterations is a classical subject, going
back to Poincare over a century ago, as described in detail in
Appendix 5. The 2D case (two real variables or one complex variable)
also goes back almost a century but the stream of literature to
which this book belongs really begins with the computer revolution
and the pioneers of scientific computation - Von Neumann, Ulam,
and so on - in the 1950s. Our subject remains an experimental
domain, and computer graphic experiments provide our main orientation.
Our fundamental tool for describing the behavior of 2D
iterations, the critical curve, was introduced by Gumowski and
Mira in 1965.2

\hfill

\textbf{1.3 PLAN OFTHE BOOK}

\hfill

In Part I, we introduce the basic concepts and vocabulary of
iteration theory, first in ID, then in 2D. We try to introduce only as
much theory as is required to understand Part 2, on exemplary
bifurcation sequences. In Part 2, we will use the vocabulary and
ideas of Part I to explain step-by-step the events in the exemplary
bifurcation sequences.

\hfill

We use the critical curves to understand the structure of attractors,
basins, basin boundaries, and their bifurcations. Then we
increase the bifurcation parameter, and explain the changes in the
configuration of attractors and basins due to bifurcations of various
types, again using the critical curves.

\hfill

These structures and changes are illustrated with still images
created by our software for the iteration of a fixed endomorphism,
based on the method of critical curves. These graphics are strung
together as movies, which may be viewed from the accompanying
CD-ROM, to give a more dynamic idea of the sequence of bifurcation
events in each of the exemplary families.

\hfill

\textbf{1.4 CONTEXT}

\hfill

Dynamics is a vast subject, and our subject is a relatively new
frontier within it. So, for those who already have an idea of the territory
of dynamical systems, we would like now to locate our
subject within this larger territory.

\hfill

Dynamical systems theory has three flavors:

\begin{itemize}
\item flows are continuous families of invertible maps generated
by a system of autonomous first-order ordinary differential
equations, and parameterized continuously by time, that is,
by real numbers;
\item cascades are discrete families of invertible maps generated
by the iteration of a given invertible map, and parameterized
discretely by the integers (zero, positive, and negative);
\item semi-cascades are discrete families of maps generated by
iteration of a given map, generally noninvertible, and
parameterized discretely by the natural numbers (zero and
the positive integers).
\end{itemize}

Both cascades and semi-cascades are also known as discrete
dynamical systems, or iterations. In this book we are primarily
interested in semi-cascades generated by a noninvertible map,
(NIM). For simplicity, we will simply call these iterations in future;
but keep in mind that all of this book belongs to the NIM flavor.

\hfill

In general, the state space, the space in which a dynamical systems
is defined, may be an arbitrary space of any dimension: 1,2,3,
and so on. This suggests a tableau of types of dynamical systems, as
shown in Figure I-I. In this tableau, there is a relationship between
cells on the same diagonal (marked with an A): In each row, the
marked cell is the cell of lowest dimension in which chaos occurs.
Hence, the tableau is called the stairway to chaos. Here chaos
means any dynamic behavior more complicated than periodic
behavior.

\hfill

In this book, we discuss only the iteration of noninvertible
maps, and the only state spaces we consider are the one-dimensional
Euclidean line and the two-dimensional Euclidean plane. In
fact, the latter is our primary subject. The ID case has been extensively
treated in recent literature (see MI) and shares the stairway
to chaos with 2D cascades and 3D flows, the contexts for the early
history of chaos theory. (See Appendix 5.) The second diagonal,
marked with B here, may be regarded as the current frontier of
chaos theory.

\hfill

\textbf{1.5 BASIC CONCEPTS OF ITERATION THEORY}

\hfill

This section introduces the basic tenninology. These concepts
will be explained in detail in ID and 2D in the next two chapters.

\hfill

Iterated map: An iterated map is the generator of a discrete
dynamical system; generally a noninvertible, continuous map.

\hfill

Multiplicity: Our maps are usually noninvertible, that is, manyto-
one, so a given point may have several preimages. The range set
may be decomposed into with zones of constant multiplicity
(bounded by critical points or critical curves) in which all points
have the same number (called the multiplicity) of preimages. These
multiple preimages detennine a tree of partial inverses for the map.
Multiplicities (explained further in the next chapter) playa fundamental
role in our theory, analogous to the degree of a polynomial
function.

\hfill

Critical sets: These sets are boundaries of zones of constant
multiplicity; thus, they separate zones of different mUltiplicity.
They consist of points with coincident inverses.

\hfill

Zones: The zones of constant multiplicity play a very fundamental
role in our view of NIM theory, analogous to the role of
degree of a polynomial in algebra.

\hfill

Partial inverses: By restricting our attention to a zone of constant
multiplicity in the range of a map, say multiplicity k, we may
define k inverses to the map. These partial inverses play the role, in
the NIM context, of the inverse of an invertible map.

\hfill

Trajectory: A trajectory embodies the basic data of a dynamical
system. It consists of the list of locations of the images of a particular
point, called the initial point, under the iterations of the map
generating the dynamical system. It is an ordered sequence (as
opposed to a set) of points.

\hfill

Attractors, basins, boundaries: These are the chief characteristic
features of an iteration, from the qualitative point of view.
Attractors are limit sets of trajectories of initial points filling an
open set, which is the basin of the attractor. The boundaries of these
basins are particularly important in applications of the theory.

\hfill

Portrait. The state space of a dynamical system may be decomposed
into a set of open sets (basins), in each of which is a single
attractor. The boundaries of these basins are particularly important
in applications of dynamical systems theory.

\hfill

Bifurcations: These are fundamental changes in the qualitative
behavior of a dynamical system, occurring as a control parameter is
varied. At certain critical values of the parameter, the qualitative
behavior of the trajectories of the system suddenly changes in a significant
way. These sudden changes, called bifurcations, usually
occur in sequences, called bifurcation sequences.

\hfill

\textbf{1.6 THE FAMILIES OF MAPS}

\hfill

The first family of maps we use to illustrate the basic ideas of
discrete dynamics is from a paper of Kawakami and Kobayashi,
studied also in a paper of Mira and coworkers:

\[
u = ax+ y
\]
\[
V = b+x^2
\]

Usually, we fix the value of a, and vary b to exhibit a bifurcation
sequence. We make use of three special cases in Part 2:

\begin{itemize}
\item Case 1: a = 0.7 (Chapter 4, Absorbing Areas)
\item Case 2: a == 1.0 (Chapter 5, Holes)
\item Case 3: a = -1.5 (Chapter 6, Fractal Boundaries).
\end{itemize}

These parameter ranges are shown in Figure 1-2.

In Chapter 7, Chaotic Contact Bifurcations, we use the double
logistic map studied by Gardini and coworkers,

\[
u = (1- c)x + 4cy(l- y)
\]
\[
v == (1 - c)y + 4cx( 1 - x)
\]

with c in [0, 1].

\hfill

\textbf{1.7 CRITICAL POINTSAND CURVES}

\hfill

The theory of critical curves for maps of the plane provides
powerful tools for locating the chief characteristic features of a discrete
dynamical system in two dimensions: the location of its
chaotic attractors, its basin boundaries, and the mechanisms of its
bifurcations. We next introduce the basic concepts of this theory
first in ID, then in 2D.

\begin{thebibliography}{9}
\bibitem{texbook}
RALPH H. ABRAHAM, LAURA GARDIINI, CHRISTIAN MIIRA, (2016) \emph{Chaos Discrete Dynamical Systems}, Springer.

\end{thebibliography}

\end{document}


